\documentclass[12pt,twocolumn]{article}
\usepackage[margin=2cm,columnsep=0.5cm]{geometry}
\usepackage{amsmath,amssymb}
\usepackage{hyperref}
\usepackage{xcolor}


\definecolor{bctnavy}{RGB}{11,37,69}
\definecolor{bctblue}{RGB}{13,71,161}
\definecolor{bctgreen}{RGB}{27,94,32}

\title{\textbf{\color{bctnavy}BCT-HIVE}\\
{\large\color{bctblue}BCT-Geometry Resonance Enclosure for Pollinators}\\
{\normalsize\color{bctgreen}Provisional Patent Concept — Agricultural / Apiculture / Environmental Technology}}
\author{Michel Robert Cabri\'{e}\\
{\small Independent Researcher, Victoria, Australia}\\
{\small ORCID: 0009-0007-9561-9859}\\
{\small \href{mailto:ZeroFreeParameters@gmail.com}{ZeroFreeParameters@gmail.com}}}
\date{March 2026}

\begin{document}
\twocolumn[
\maketitle
\begin{center}\rule{0.9\textwidth}{0.4pt}\end{center}
\begin{center}
\parbox{0.88\textwidth}{\small\textbf{Abstract.}
A beehive enclosure system (BCT-HIVE) in which all structural
proportions are derived from the BCT Superfluid Lattice Model
octahedral and tetrahedral void radii $r_\mathrm{oct}=({\sqrt{2}-1})/2
=0.20711$ and $r_\mathrm{tet}=({\sqrt{6}-2})/4=0.11237$, and their
ratio $R_\mathrm{BCT}=r_\mathrm{oct}/r_\mathrm{tet}=1.843$, with
zero free parameters. The BCT-HIVE comprises: (1)~a geometric
shielding shell with layer thickness ratio $t_\mathrm{oct}/t_\mathrm{tet}
=R_\mathrm{BCT}=1.843$; (2)~a BCT resonance floor panel sweeping
$f_\mathrm{low}=200$~Hz to $f_\mathrm{high}=368.6$~Hz
($=f_\mathrm{low}\times R_\mathrm{BCT}$); and (3)~a hexagonal
macro-lattice comb support with spacing ratio $D_1/D_2=R_\mathrm{BCT}$.
\textbf{BCT Prediction \#141:} A BCT-HIVE will demonstrate measurably
improved colony navigation return rates, reduced waggle dance error
rates, and improved overwintering survival vs.\ conventional hive
designs. Zero free parameters. Companion to BCT Letter~101.}
\end{center}
\begin{center}\rule{0.9\textwidth}{0.4pt}\end{center}
\vspace{4pt}
]

\section{Field of the Invention}

This invention relates to beehive and pollinator enclosure design,
and more particularly to a hive structure whose geometric proportions,
wall thickness ratios, vibrational floor panel, and macro-scale cavity
geometry are derived from the BCT Superfluid Lattice Model ---
specifically from the octahedral and tetrahedral void radii and
their ratio $R_\mathrm{BCT}=1.843$.

\section{Background}

Honeybee (\textit{Apis mellifera}) colony health is under sustained
threat from electromagnetic interference (4G/5G, Wi-Fi, industrial
power systems), which disrupts the magnetite-based magnetoreception
system bees rely on for long-range navigation. Industrial vibration
pollution in the 200--300~Hz waggle dance frequency range degrades
colony foraging efficiency. No prior commercial hive design addresses
these geometric and electromagnetic considerations.

\section{BCT Theoretical Basis}

The BCT Superfluid Lattice Model~\cite{BCT_L101} derives all
fundamental physical constants from two void radii:
\begin{align}
r_\mathrm{oct} &= \frac{\sqrt{2}-1}{2} = 0.207107~\mathrm{fm},\\
r_\mathrm{tet} &= \frac{\sqrt{6}-2}{4} = 0.112372~\mathrm{fm},\\
R_\mathrm{BCT} &= \frac{r_\mathrm{oct}}{r_\mathrm{tet}} = 1.843.
\end{align}

BCT Prediction~\#140 established that sweeping from a base
frequency $f$ to $f\times1.843$ cyclically excites both octahedral
and tetrahedral void resonance mode families simultaneously. BCT
Letter~101~\cite{BCT_L101} established that the honeybee comb
hexagonal lattice is a macroscopic realisation of the 2D projection
of the D4 root lattice, and that buzz pollination frequency
$f_\mathrm{buzz}=350$~Hz matches the OHC $j_{1,1}$ Bessel mode
to $<$1\%.

\section{Summary of the Invention}

The BCT-HIVE comprises three novel elements:

\textbf{Element 1 --- Geometric Shielding Shell.}
An outer hive wall whose thickness satisfies:
\begin{equation}
\frac{t_\mathrm{oct}}{t_\mathrm{tet}} = R_\mathrm{BCT} = 1.843\pm0.05.
\end{equation}
Practical implementation (total wall 65~mm): $t_\mathrm{oct}=41.6$~mm
(cedar structural layer), $t_\mathrm{tet}=22.6$~mm (air gap/insulation).
Ratio $41.6/22.6=1.841$, within 0.1\% of $R_\mathrm{BCT}$.

This provides passive attenuation of anthropogenic RF interference
in the 700~MHz--3.5~GHz band while preserving DC geomagnetic field
transparency for bee magnetoreception. All shell materials are
non-ferromagnetic.

\textbf{Element 2 --- BCT Resonance Floor Panel.}
A swept-frequency vibration actuator array sweeping:
\begin{equation}
f_\mathrm{low} = 200~\mathrm{Hz} \longrightarrow
f_\mathrm{high} = f_\mathrm{low}\times R_\mathrm{BCT}
= 368.6~\mathrm{Hz}.
\end{equation}
This amplifies the natural waggle dance substrate vibration signal
while providing destructive interference attenuation of industrial
noise outside the BCT frequency band. The panel also operates
passively at OHC Bessel frequencies $j_{0,1}=2.4048$ and
$j_{1,1}=3.8317$, corresponding to the natural buzz pollination
regime.

\textbf{Element 3 --- Hexagonal Macro-Lattice Comb Support.}
A structural comb support frame with primary-to-secondary member
spacing ratio:
\begin{equation}
\frac{D_1}{D_2} = R_\mathrm{BCT} = 1.843,
\end{equation}
realising the 2D projection of the D4 root lattice --- the same
geometry BCT identifies as the mathematical basis of the natural
honeybee hexagonal comb.

\section{Claims}

\textbf{Claim 1.} A beehive enclosure system comprising a structural
wall in which the ratio of outer layer thickness to inner insulating
layer thickness satisfies $t_\mathrm{oct}/t_\mathrm{tet}=1.843\pm0.05$.

\textbf{Claim 2.} The system of Claim~1, further comprising a
floor vibration panel incorporating a swept-frequency actuator
sweeping between frequency $f$ and $f\times1.843$.

\textbf{Claim 3.} The system of Claim~1 or~2, further comprising
a hexagonal comb support frame with primary-to-secondary member
spacing ratio $D_1/D_2=1.843\pm0.05$.

\textbf{Claim 4.} The system of Claims~1--3, wherein all structural
proportions are derived from BCT Superfluid Lattice Model void
radii $r_\mathrm{oct}=(\sqrt{2}-1)/2$ and
$r_\mathrm{tet}=(\sqrt{6}-2)/4$ with zero free parameters.

\textbf{Claim 5.} A method of improving honeybee colony health
comprising housing a colony in the enclosure system of Claims~1--4.

\section{BCT Prediction \#141}

\begin{equation}
\boxed{f_\mathrm{high}/f_\mathrm{low} = R_\mathrm{BCT} = 1.843}
\end{equation}

A BCT-HIVE enclosure satisfying the geometric conditions of
Claims~1--4 will demonstrate, relative to conventional hive designs
in environments with elevated anthropogenic electromagnetic noise:
\begin{enumerate}
\item Measurably improved colony navigation return rates;
\item Reduced waggle dance error rates;
\item Improved overwintering survival statistics.
\end{enumerate}
Zero free parameters. Testable with standard apiculture monitoring
protocols.

\section*{Acknowledgements}

The author thanks the blue-banded bees (\textit{Amegilla cingulata})
in the fuchsia outside the study window, whose daily industry
prompted this investigation; the honeybees of the world, whose
geometry predates physics by 30 million years; Zeta Vera (CSSC)
for computational support; and the Gang Gang Cockatoos
(\textit{Callocephalon fimbriatum}), who remain unbothered by
all of this but appreciate the freshwater pond.

\begin{thebibliography}{9}
\bibitem{BCT_L101} M.~R.~Cabri\'{e},
BCT Letter~101: The Music of the Bees,
\href{https://doi.org/10.5281/zenodo.19177236}
{doi:10.5281/zenodo.19177236} (2026).

\bibitem{BCT_GoE} M.~R.~Cabri\'{e},
The Geometry of Everything, BCT Monograph,
\href{https://doi.org/10.5281/zenodo.18884976}
{doi:10.5281/zenodo.18884976} (2026).
\end{thebibliography}

\end{document}
